\section{Countermodels and kill criteria}
\label{sec:countermodels}

The following examples prevent the regular-kernel theorem from being
silently promoted to the literal carrier.

\subsection{Bounded amplitude is insufficient}

Define on \([1,U]\times[1,V]\)
\[
 K_{\rm match}(u,v)=\mu(uv).
\]
Then \(\|K_{\rm match}\|_\infty\le1\), but
\begin{equation}
 \cS(K_{\rm match})
 =
 \sum_{u\le U,v\le V}\mu(uv)^2
 =
 \mathfrak S_{\rm sf}UV+o(UV).
\label{eq:sign-matched-countermodel}
\end{equation}
Thus no estimate of the form \(o(UV)\) can hold for all bounded kernels.
The kernel has selected the exact arithmetic sign.

There is also a nonnegative version:
\[
 K_+(u,v)=\one_{\mu(uv)=1}.
\]
By \eqref{eq:squarefree-pair-asymptotic} and
\eqref{eq:primitive-rectangle},
\begin{equation}
 \cS(K_+)
 =
 \frac{\mathfrak S_{\rm sf}}2UV+o(UV).
\label{eq:nonnegative-sign-selected-mask}
\end{equation}
Hence nonnegativity of the mask does not prevent coherent parity selection.
These are abstract plane kernels; we do not assert that the physical
prime-line mask equals either one.

\subsection{A prime mask can freeze the plane sign}

Let
\[
 K_{\mathbb P,N}(u,v)
 =
 \one_{\substack{u,v\le N\ {\rm prime}\\u\ne v}}.
\]
Its support is primitive and squarefree.  Since \(\mu(p)=-1\) for every
prime,
\[
 \cS(K_{\mathbb P,N})
 =
 \pi(N)(\pi(N)-1)
 =
 \|K_{\mathbb P,N}\|_{\ell^1}.
\]
Thus prime support by itself does not create M\"obius cancellation; it can
freeze the sign completely.  This example is not the literal two-prime-line
kernel, but it explains why an upper-bound sieve for its capacity cannot
serve as a variation or discrepancy theorem.

\subsection{A primitive line contains a Chowla-type boundary}

Take
\[
 K_{\rm line,N}(u,v)=\one_{v=u+1,\ u\le N}.
\]
Every supporting pair is coprime, but
\[
 \cS(K_{\rm line,N})
 =
 \sum_{u\le N}\mu(u)\mu(u+1).
\]
An unweighted \(o(N)\) theorem here is the open two-point Chowla
correlation.  The regular-kernel theorem correctly does not apply: a thin
diagonal graph has mixed variation growing like its line length, rather
than \(O(\|K\|_\infty)\).  Known logarithmically averaged results
\citep{Tao2016LogChowla} do not turn this raw fixed line into an
unweighted theorem.

\subsection{Small ambient bound can be physically vacuous}

Let \(K\) be supported on a set of size \(A=o(UV)\), with
\(|K|\le1\).  The unconditional conclusion
\[
 \cS(K)=o(UV)
\]
does not imply \(o(A)\), and may be larger than the entire physical mass.
The only relevant statement is the relative certificate
\[
 \cB_D(K)=o(\cA(K))
\]
with \(\cA(K)\) chosen from the actual normalization.  This is why
\cref{def:variation-tail-gate} is not cosmetic.

\subsection{Pairwise decay can lose to row growth}

For each dimension \(R\), let a symmetric defect have zero diagonal and
all off-diagonal entries \(R^{-1}\).  Every fixed entry tends to zero, but
the top eigenvalue is
\[
 \frac{R-1}{R}\longrightarrow1.
\]
Therefore a fixed-pair primitive-plane estimate, even if uniform in all
other parameters, is not an operator saving unless its row sum or Perron
radius also tends to zero.

\subsection{A mixed trace does not control the erasure operator}

Let
\[
 R=0,
 \qquad
 E=
 \begin{pmatrix}
 0&1\\
 1&0
 \end{pmatrix}.
\]
Then every \(ER\) rectangle scalar and
\(\operatorname{tr}(ER)\) vanish, while
\[
 \|E\|=1.
\]
Thus even perfect cancellation in the anchored TPC--70 target cannot be
renamed as an erasure-operator bound.  A direct two-corner or complete
moment bridge is indispensable.

\subsection{Separate \texorpdfstring{\(q\)}{q}-bounds can destroy cancellation}

Take two divisor slices with
\[
 K_{q_1}=F,\qquad K_{q_2}=-F.
\]
The aggregate kernel is zero, while
\[
 \cV_\square(K_{q_1})+\cV_\square(K_{q_2})
 =
 2\cV_\square(F).
\]
Thus the q-resolved triangle inequality can be strictly weaker than
variation of the exact aggregate.  Conversely, no cancellation between
q-slices may be assumed unless their literal phases produce it.

\subsection{Failure of one certificate is not failure of cancellation}

Let \(K\) be any highly oscillatory kernel whose signed plane sum happens
to vanish exactly.  Its mixed variation can still be of order \(UV\).
Therefore a large value of
\(\cV_\square(K^{[d]})\) kills the divisor--variation route only.  It does
not prove that the original physical coefficient is large.

\subsection{Explicit stop rules}

For the current route, the following outcomes are decisive.

\begin{enumerate}[label=\textup{K\arabic*.},leftmargin=2.4em]
 \item If the exact physical kernel has
 \[
 \cB_D(K^\circ_{\mathfrak a})
 \not=o(\cA_{\mathfrak a})
 \]
 for an anchor family carrying a nonnegligible part of the declared trace
 scale, then the present anchored variation route does not close.
 \item Even if the anchored sums make
 \(\operatorname{tr}(ER)=o(1)\) at its declared scale, the operator route
 stops unless a new direct-entry, row-local, or growing-moment bridge is
 proved.  Global trace cancellation alone is not such a bridge.
 \item If the native defect becomes small but the mass comparison,
 represented lower frame, or nuisance shorting term fails, reassembly
 stops; the plane result is not upgraded to a prime-pair statement.
 \item If a required net exponent loss reaches or exceeds \(1/400\),
 including equality, the current endpoint implication stops.
\end{enumerate}

Each stop rule leaves a valid negative or boundary result: it identifies
which exact certificate fails without claiming that every possible route
fails.
